
\section{流动稳定性}

物理上，只有稳定的流动状态才容易被观察到，即使是流体力学方程的精确解描述的流动状态，如果不稳定，也无法观察。

流动稳定行研究内容：研究不同流动状态对扰动的响应特性。具体：流动是不是稳定的，在什么条件下流动由稳定变为不稳定，不稳定流动将向什么方向发展。

\subsection{基本概念}
\subsubsection{稳定性定义}
一般研究层流稳定性，完全发展的湍流中某些现象可以用流动稳定性理论解释。

流动稳定性研究对象称为基本流（场），是一个层流流场，理论上这个流场满足NS方程组及其定解条件。

基本流场的获取是困难的，实际应用中往往是NS方程精确解的近似，如平行流假设的基本流、实验拟合的速度型等。基本流是流体动力学控制方程（如Navier-Stokes方程）在特定理想化条件（如特定的对称性、定常性、平行流等）下所对应的精确解。它作为研究其稳定性的未受扰动的背景状态。

流动稳定性的物理定义：基本流是否可以观察到。若初始时刻流动偏离了基本流，这种偏离称为对基本流的扰动。
\begin{itemize}
    \item 相空间整体稳定（globally stable in phase space）：对任何扰动均稳定
    \item 条件稳定（conditionally stable）：仅对幅值小于某一阈值的扰动稳定
    \item 渐近稳定（asymptotically stable）：扰动衰减，基本流得以恢复
        \begin{itemize}
            \item 单调稳定：扰动所有时刻衰减
            \item 暂态增长（稳定）：某个时间扰动增长，最终还是会衰减。
        \end{itemize}
    \item 中性稳定：扰动幅值大小保持不变
    \item 不稳定的：扰动被放大
\end{itemize}
不稳定的基本流不在能被观察到，流场最终演变为湍流或另一种层流，相比于基本流，失稳后的新流场，对称性一般降低。

（微分方程解的稳定性理论，可以数学地描述上面的各种情况）

线性理论里，初始扰动被限制为无穷小，实际中用充分小的扰动代替无穷小的扰动。充分小的涵义与具体流动条件有关，解析上通过忽略扰动方程中的非线性项将问题线性化，只要扰动小到忽略非线性项是合理的，这个扰动就是“充分小”的。

\subsubsection{临界参数}
某些流动参数决定了流动的稳定性特性。在参数组成的坐标系中，流动的稳定性可以被划分为稳定和不稳定区域。区域边界称为边缘稳定（marginally stable），边缘稳定一定中性稳定，反之不然，如稳定区域中可能出现孤立的中性稳定区域。

\subsection{不稳定机理}
平衡态1可以由扰动恢复，则流动稳定，如果必须改变流场建立平衡态2，则流动不稳定。流体所受作用力由动量方程描述，平衡状态就是各种力处于平衡，当各种因素影响到作用力，间接地影响基本流的稳定性特性。

外力对流动稳定性影响：浮力不稳定性、离心不稳定性

没有外力和黏性力作用时，流体根据惯性力和压力的平衡进行运动。当流体沿压力增加的方向运动会增强不稳定性。
黏性对流动稳定性具有两重性：一方面，边界限制扰动发展，起稳定化作用；另一方面，边界产生强速度剪切，通过黏性扩散导致不稳定性增强。

根据机理对流动不稳定性分类，最重要最常见的两类
\begin{itemize}
    \item Rayleigh-Taylor不稳定性：浮力、热、离心不稳定性的统称
    \item Kelvin-Helmholtz不稳定性：由速度剪切引起的不稳定性
\end{itemize}

\subsection{典型不稳定性现象}

